\documentclass[a4paper,12pt]{article}

% --- Pakiety ---
\usepackage[polish]{babel}
\usepackage[utf8]{inputenc}
\usepackage[T1]{fontenc}
\usepackage[top=2.5cm,bottom=2.5cm,left=3cm,right=3cm]{geometry}
\usepackage{amsmath}
\usepackage{graphicx}
\usepackage{booktabs}
\usepackage[colorlinks=true, allcolors=blue]{hyperref}
\usepackage{amssymb}
\usepackage{listings}
\usepackage{titlesec}

\lstset{
    basicstyle=\ttfamily\small,
    breaklines=true,
    columns=fullflexible
}

\titlespacing*{\subsection}{0pt}{12pt}{6pt}
\setlength{\parindent}{15pt}
\setlength{\parskip}{6pt}

% --- Tytuł ---
\title{
    \vspace{1.5cm}
    \textbf{PODSTAWY TEORII INFORMACJI}\\
    \vspace{1.5cm}
    Laboratorium 3\\
    Syntaktyczna analiza danych\\
    \vspace{1.5cm}
    \huge \textbf{Ekstrakcja i redukcja informacji w danych meteorologicznych}\\
    \vspace{1cm}
}

\author{
    Maciek Chotkowski \\ Martsin Lazouski \\ Nadia Ziecina \\
    Michał Kasznia \\ Julian Kasiński \\ Volodymyr Kamenchuk \\ Usevalad Nikalaichuk 
}

\date{\today}

\begin{document}

\maketitle
\newpage

\tableofcontents
\newpage

%--------------------------------------------------

\section{Wprowadzenie}

W poprzednim etapie projektu konieczne było sprowadzenie analizowanego problemu do uproszczonej postaci, umożliwiającej dalsze przetwarzanie danych. Ograniczono się do wybranych parametrów meteorologicznych oraz ich reprezentacji w postaci ciągów liczbowych.

\noindent
Dopiero w niniejszym etapie wprowadzamy konkretne rozwiązanie obliczeniowe. Oznacza to przejście od opisu problemu do jego realizacji w postaci narzędzia przetwarzającego dane.

\noindent
Celem ćwiczenia jest analiza syntaktyczna danych, polegająca na ich przetwarzaniu i ocenie przy użyciu formalnych kryteriów obliczeniowych.

%--------------------------------------------------

\section{Etap 1: Narzędzie pozyskiwania danych}

Pierwszym etapem opracowanego systemu jest moduł automatycznego pobierania
danych meteorologicznych i zapisania ich w prostej strukturze plików.
Na tym etapie nie tworzymy jeszcze rankingu ani jednego wspólnego pliku
porównawczego. Celem jest jedynie zebranie danych i rozdzielenie ich według
źródła oraz typu wielkości fizycznej.

\noindent
Każde źródło danych otrzymuje własny katalog. W katalogu znajdują się trzy
pliki CSV:
\begin{verbatim}
nazwa_zrodla/
    temp.csv
    opady.csv
    wiatr.csv
\end{verbatim}

\noindent
Taki zapis jest zgodny z poziomem syntaktycznym analizy, ponieważ każdy plik
zawiera prosty ciąg liczbowy jednej cechy. Temperatura, opady i wiatr nie są
jeszcze interpretowane semantycznie. Są traktowane jako sygnały liczbowe,
które można później dopasować po czasie i porównać z innymi źródłami.

\noindent
Wejściem do modułu są:
\begin{itemize}
    \item współrzędne geograficzne,
    \item horyzont prognozy,
    \item nazwa źródła danych.
\end{itemize}

\noindent
Wyjściem są trzy uporządkowane sygnały:
\[
T = \{T_1, T_2, \dots, T_n\}, \quad
P = \{P_1, P_2, \dots, P_n\}, \quad
W = \{W_1, W_2, \dots, W_n\}
\]

\noindent
Dane są zapisywane osobno dla każdego źródła i przekazywane do dalszych
etapów analizy.
\newpage
\subsection{Implementacja narzędzia}

Poniższy fragment przedstawia podstawową wersję programu pozyskującego dane
z dwóch źródeł: Open-Meteo oraz yr.no. Program pobiera z odpowiedzi API tylko
pola potrzebne w projekcie: czas, temperaturę, opady i prędkość wiatru.
Następnie zapisuje każdą wielkość do osobnego pliku w katalogu danego
źródła.

\noindent
Wynikiem działania kodu nie jest jeden plik zbiorczy. Wynikiem jest katalog
\texttt{open\_meteo} oraz katalog \texttt{yr\_no}. W każdym z nich znajdują
się pliki \texttt{temp.csv}, \texttt{opady.csv} oraz \texttt{wiatr.csv}.
Dzięki temu po etapie zbierania danych wiadomo, z którego źródła pochodzi
każdy sygnał i jakiego parametru dotyczy.

\begin{lstlisting}[language=Python]
import requests, pandas as pd
from pathlib import Path

LAT, LON = 52.23, 21.01

om = requests.get(
    "https://api.open-meteo.com/v1/forecast",
    params={
        "latitude": LAT,
        "longitude": LON,
        "hourly": "temperature_2m,precipitation,wind_speed_10m",
        "forecast_days": 3
    }
).json()

df_om = pd.DataFrame({
    "time": om["hourly"]["time"],
    "temp": om["hourly"]["temperature_2m"],
    "opady": om["hourly"]["precipitation"],
    "wiatr": om["hourly"]["wind_speed_10m"]
})

yr = requests.get(
    "https://api.met.no/weatherapi/locationforecast/2.0/compact",
    headers={"User-Agent": "pti-lab"},
    params={"lat": LAT, "lon": LON}
).json()

df_yr = pd.DataFrame([{
    "time": item["time"][:16],
    "temp": item["data"]["instant"]["details"]["air_temperature"],
    "opady": item["data"].get("next_1_hours", {}).get("details", {}).get("precipitation_amount", 0),
    "wiatr": item["data"]["instant"]["details"]["wind_speed"]
} for item in yr["properties"]["timeseries"]])

def zapisz_zrodlo(nazwa, df):
    out = Path(nazwa)
    out.mkdir(exist_ok=True)
    df[["time", "temp"]].to_csv(out / "temp.csv", index=False)
    df[["time", "opady"]].to_csv(out / "opady.csv", index=False)
    df[["time", "wiatr"]].to_csv(out / "wiatr.csv", index=False)

zapisz_zrodlo("open_meteo", df_om)
zapisz_zrodlo("yr_no", df_yr)
\end{lstlisting}

\subsection{Zakres i postać pobieranych danych}

W ramach narzędzia pobierane są prognozowane dane meteorologiczne
dla lokalizacji Warszawa (52.23N, 21.01E).

\noindent
Uwzględniono następujące parametry:
\begin{itemize}
    \item temperatura powietrza (°C),
    \item opady (mm),
    \item prędkość wiatru (m/s).
\end{itemize}

\noindent
Dane pobierane są z API w formacie JSON, który ma strukturę
hierarchiczną (zagnieżdżone pola i listy). Poszczególne wartości
odczytywane są z odpowiednich pól, np.:

\begin{verbatim}
Open-Meteo: hourly -> temperature_2m
yr.no: properties -> timeseries -> instant -> details
\end{verbatim}

\noindent
Po ekstrakcji dane przekształcane są do postaci tabelarycznej
oraz wektorowej:

\[
T(t), \quad P(t), \quad W(t)
\]

\noindent
gdzie:
\begin{itemize}
    \item $T(t)$ – temperatura,
    \item $P(t)$ – opady,
    \item $W(t)$ – prędkość wiatru.
\end{itemize}

\noindent
Ostatecznie dane nie są jeszcze łączone w jeden plik. Każdy parametr jest
zapisywany osobno w katalogu źródła:
\begin{verbatim}
open_meteo/temp.csv
open_meteo/opady.csv
open_meteo/wiatr.csv
yr_no/temp.csv
yr_no/opady.csv
yr_no/wiatr.csv
\end{verbatim}
Analogicznie można zapisać kolejne źródła, np. \texttt{graphcast}
albo \texttt{stacja\_pomiarowa}. Dopiero w następnym
etapie pliki te są dopasowywane po czasie i porównywane.


\subsection{Rozszerzenie: pobieranie danych z modeli AI dostępnych w Open-Meteo}

Oprócz klasycznych serwisów pogodowych wykorzystano również prognozy
generowane przez modele sztucznej inteligencji dostępne w API Open-Meteo.
W projekcie użyto modeli GraphCast oraz AIFS, które traktowane są jako
dodatkowe źródła danych meteorologicznych.

\noindent
Podobnie jak w przypadku Open-Meteo i yr.no, dane pobierane są w postaci
szeregów czasowych dla wybranej lokalizacji. Końcowym wynikiem nie jest
globalna prognoza atmosferyczna, lecz proste wektory wartości:
\[
temp = [temp_1, temp_2, \dots, temp_n]
\]
\[
precip = [precip_1, precip_2, \dots, precip_n]
\]
\[
wind = [wind_1, wind_2, \dots, wind_n]
\]

\noindent
Modele AI są udostępniane przez Open-Meteo za pomocą parametru
\texttt{models}. Dzięki temu można pobierać prognozy różnych modeli przy
użyciu tego samego endpointu API. W projekcie wykorzystano:
\begin{itemize}
    \item \texttt{gfs\_graphcast025} -- model GraphCast,
    \item \texttt{ecmwf\_aifs025\_single} -- model AIFS.
\end{itemize}

\noindent
Dane pobierane są bezpośrednio z API w formacie JSON, a następnie
przekształcane do postaci tabelarycznej przy użyciu biblioteki
\texttt{pandas}. Każde źródło otrzymuje osobny katalog zawierający trzy
pliki CSV:
\begin{verbatim}
graphcast/
    temp.csv
    opady.csv
    wiatr.csv

aifs/
    temp.csv
    opady.csv
    wiatr.csv
\end{verbatim}

\noindent
Wszystkie źródła posiadają identyczną strukturę danych, dzięki czemu
mogą być później analizowane tym samym kodem MATLAB lub Python.
\newpage
\subsubsection{Postać danych AI}

Modele AI zwracają prognozy dla określonego punktu geograficznego oraz
wybranego przedziału czasu. W projekcie wykorzystywane są następujące
parametry meteorologiczne:
\begin{itemize}
    \item temperatura powietrza (\texttt{temperature\_2m}),
    \item opady (\texttt{precipitation}),
    \item prędkość wiatru (\texttt{wind\_speed\_10m}).
\end{itemize}

\noindent
Po pobraniu dane są zapisywane do tabeli:
\begin{verbatim}
time, temp, opady, wiatr
\end{verbatim}

\noindent
Każda kolumna reprezentuje jeden sygnał czasowy wykorzystywany później
w analizie błędów oraz porównaniu modeli.

\subsubsection{Przykładowa implementacja pobierania danych AI}

Poniższy kod przedstawia przykład pobierania prognoz z modeli AI
udostępnianych przez Open-Meteo. Program pobiera dane dla Warszawy,
a następnie zapisuje temperaturę, opady oraz wiatr do osobnych plików CSV.

\begin{lstlisting}[language=Python]
import requests
import pandas as pd
from pathlib import Path

LAT = 52.23
LON = 21.01

DATA_START = "2026-05-10"
DATA_END = "2026-05-13"

MODELS = {
    "graphcast": "gfs_graphcast025",
    "aifs": "ecmwf_aifs025_single"
}

for source_name, model_name in MODELS.items():

    url = "https://api.open-meteo.com/v1/forecast"

    params = {
        "latitude": LAT,
        "longitude": LON,
        "start_date": DATA_START,
        "end_date": DATA_END,
        "hourly": (
            "temperature_2m,"
            "precipitation,"
            "wind_speed_10m"
        ),
        "wind_speed_unit": "ms"
    }

    response = requests.get(
        url,
        params=params
    )

    data = response.json()

    if "hourly" not in data:

        print(f"Błąd dla {source_name}")
        print(data)
        continue

    df = pd.DataFrame({
        "time": data["hourly"]["time"],
        "temp": data["hourly"]["temperature_2m"],
        "opady": data["hourly"]["precipitation"],
        "wiatr": data["hourly"]["wind_speed_10m"]
    })

    out = Path(source_name)
    out.mkdir(exist_ok=True)

    df[["time", "temp"]].to_csv(
        out / "temp.csv",
        index=False
    )

    df[["time", "opady"]].to_csv(
        out / "opady.csv",
        index=False
    )

    df[["time", "wiatr"]].to_csv(
        out / "wiatr.csv",
        index=False
    )

    print(f"Zapisano dane dla: {source_name}")
\end{lstlisting}
\newpage

\noindent
Po zapisaniu danych każdy model AI posiada identyczny układ plików,
co pozwala później połączyć wszystkie źródła w jedną strukturę
analityczną.

\noindent
Tak przygotowane dane mogą zostać następnie wczytane do środowiska MATLAB
lub Python i połączone z danymi rzeczywistymi oraz prognozami klasycznymi.
Ujednolicona struktura plików pozwala wykonywać dalszą analizę przy użyciu
tych samych procedur obliczeniowych niezależnie od źródła prognozy.

\noindent
W praktyce oznacza to, że modele AI są traktowane identycznie jak pozostałe
źródła danych meteorologicznych. Możliwe staje się więc bezpośrednie
porównywanie temperatury, opadów oraz prędkości wiatru pomiędzy różnymi
modelami oraz obliczanie miar błędu, takich jak RMSE.

\subsubsection{Rozszerzanie systemu o kolejne modele AI}

Dodanie kolejnego modelu AI polega wyłącznie na dopisaniu nowego wpisu
do słownika \texttt{MODELS}. Dzięki temu nie ma potrzeby zmiany dalszego
kodu odpowiedzialnego za analizę danych.

\noindent
Przykładowo:
\begin{verbatim}
MODELS = {
    "graphcast": "gfs_graphcast025",
    "aifs": "ecmwf_aifs025_single",
    "aigfs": "ncep_aigfs025"
}
\end{verbatim}

\noindent
Wpływ na dalszy etap projektu pozostaje więc niewielki, ponieważ każde
źródło zapisuje dane w tym samym formacie CSV. Dzięki temu obliczanie
RMSE, filtrowanie danych oraz tworzenie wykresów działa identycznie dla
modeli klasycznych i modeli AI.

\newpage


%--------------------------------------------------
\section{Etap 2: Przetwarzanie, analiza i ewaluacja danych}

\subsection{Automatyzacja procesu analizy}

W drugim etapie projektu opracowano zautomatyzowany system analizy jakości prognoz meteorologicznych generowanych przez różne źródła danych. Implementację wykonano w języku Python z wykorzystaniem bibliotek \texttt{pandas}, \texttt{numpy} oraz \texttt{matplotlib}. Głównym celem tego etapu było stworzenie jednolitego środowiska umożliwiającego jednoczesne porównanie modeli AIFS, GraphCast, Open-Meteo, yr.no oraz klasycznego modelu GFS względem danych referencyjnych (\textit{ground truth}). 

\noindent
W pierwszej kolejności zaimplementowano procedurę automatycznego wczytywania danych z plików CSV. Każdy model posiadał własny katalog zawierający szeregi czasowe temperatury. Dane referencyjne zostały zdefiniowane jako osobne źródło \texttt{default\_model}, względem którego porównywano wszystkie pozostałe prognozy. Program przed rozpoczęciem analizy weryfikował poprawność struktury danych oraz obecność wymaganych kolumn \texttt{time} i \texttt{temp}. W przypadku braku pliku lub niepoprawnego formatu dane były automatycznie odrzucane z dalszej analizy. 

\noindent
Kolejnym krokiem była synchronizacja szeregów czasowych. Wszystkie znaczniki czasu zostały skonwertowane do wspólnego formatu typu \texttt{datetime}, a następnie dane sortowano chronologicznie oraz usuwano duplikaty. Po przygotowaniu danych wykonywano operację \texttt{merge}, która łączyła prognozy modeli z danymi referencyjnymi wyłącznie dla wspólnych godzin pomiarowych. Takie podejście gwarantowało, że porównanie odbywało się dla identycznych momentów czasowych.

\subsection{Metody oceny jakości prognoz}

Ocena jakości prognoz została oparta na dwóch równoległych mechanizmach ewaluacyjnych. Pierwszym z nich był błąd bezwzględny, obliczany jako wartość bezwzględna różnicy pomiędzy temperaturą referencyjną a temperaturą prognozowaną przez dany model:

\[
|T_{ref} - T_{model}|
\]

\noindent
Na podstawie błędu bezwzględnego zdefiniowano autorski system punktacji rankingowej. Dla każdej godziny prognozy model otrzymywał określoną liczbę punktów zależną od wielkości odchylenia od wartości referencyjnej:

\begin{itemize}
    \item 3 punkty -- błąd mniejszy lub równy 0.5°C,
    \item 2 punkty -- błąd mniejszy lub równy 1.5°C,
    \item 1 punkt -- błąd mniejszy lub równy 3.0°C,
    \item 0 punktów -- błąd większy niż 3.0°C.
\end{itemize}

\noindent
Takie podejście umożliwia ocenę nie tylko średniej jakości modelu, ale również jego stabilności oraz zdolności do utrzymywania niewielkich odchyleń przez dłuższy okres czasu. Implementacja systemu punktowego została zrealizowana jako osobna funkcja przypisująca punkty na podstawie wielkości błędu. 

\noindent
Drugim elementem ewaluacji była klasyczna metryka RMSE (\textit{Root Mean Square Error}), czyli pierwiastek błędu średniokwadratowego:

\[
RMSE = \sqrt{\frac{1}{n}\sum_{i=1}^{n}(y_i - \hat{y}_i)^2}
\]

\noindent
Metryka ta silniej penalizuje duże odchylenia prognozy od wartości rzeczywistych, dzięki czemu pozwala wykrywać modele generujące sporadyczne, lecz bardzo duże błędy. Wartość RMSE obliczano dla całego wspólnego zakresu czasowego pomiędzy modelem a danymi referencyjnymi. 

\noindent
Na podstawie uzyskanych wyników tworzony był końcowy ranking modeli. System automatycznie sortował źródła według procentowej skuteczności punktowej, a w przypadku zbliżonych rezultatów dodatkowym kryterium była wartość RMSE. Dla każdego modelu prezentowano liczbę zdobytych punktów, maksymalną możliwą liczbę punktów, procent skuteczności, wartość RMSE oraz liczbę wspólnych godzin uwzględnionych w analizie. 

\subsection{Wizualizacja i interpretacja wyników}

Istotnym elementem analizy była również wizualna weryfikacja jakości prognoz. W tym celu opracowano moduł generujący wykres porównawczy przedstawiający przebieg temperatury referencyjnej oraz prognoz wszystkich analizowanych modeli. Dane zostały przedstawione na wspólnym układzie współrzędnych z wykorzystaniem różnych kolorów i stylów linii dla każdego źródła. Krzywa danych referencyjnych została dodatkowo wyróżniona większą grubością linii, co ułatwia ocenę zgodności poszczególnych modeli z rzeczywistym przebiegiem temperatury. 

\noindent
Analiza wizualna pozwala na identyfikację charakterystycznych zachowań modeli, których nie zawsze można jednoznacznie uchwycić wyłącznie przy pomocy pojedynczych metryk liczbowych. Dzięki wykresom możliwe jest określenie, czy dany model poprawnie odwzorowuje dynamikę zmian temperatury, lokalne maksima i minima oraz amplitudę wahań dobowych. Pozwala to również zauważyć zjawiska takie jak wygładzanie sygnału, przesunięcia fazowe, generowanie fałszywych pików czy systematyczne przeszacowanie lub niedoszacowanie temperatury.

\noindent
Wygenerowane wykresy zapisywano automatycznie do pliku \texttt{wykres\_temperatura.png}, co umożliwia ich bezpośrednie wykorzystanie w dalszej części raportu oraz podczas analizy porównawczej poszczególnych modeli prognostycznych. 

\subsection{Analiza z wykorzystaniem danych rzeczywistych}

Dane prognozowane zostały zestawione z rzeczywistymi pomiarami
meteorologicznymi zapisanymi w pliku CSV.

\noindent
Po dopasowaniu danych względem czasu przeprowadzono analizę
różnic pomiędzy prognozą a rzeczywistością.

\noindent
Na tej podstawie obliczono błąd RMSE, który stanowi miarę
dokładności prognozy.


%--------------------------------------------------
\subsection{Wyniki}

\subsubsection{Porównywanie modeli na podstawie przygotowanych plików CSV}

Właściwe porównanie zostało wykonane po etapie pobierania danych. Wszystkie
źródła zostały sprowadzone do tej samej struktury katalogów:
\begin{verbatim}
zrodlo/
    temp.csv
    opady.csv
    wiatr.csv
\end{verbatim}

\noindent
Jako punkt odniesienia przyjęto katalog \texttt{default\_model}. W naszym
programie pełni on rolę danych referencyjnych, czyli tzw. \textit{ground truth}.
Pozostałe źródła, czyli \texttt{open\_meteo}, \texttt{yr\_no},
\texttt{gfs}, \texttt{aifs} oraz \texttt{graphcast}, są porównywane z tym
samym punktem odniesienia.

\noindent
Przed obliczeniem błędu dane są łączone po kolumnie \texttt{time}. Oznacza to,
że porównywane są tylko wartości odnoszące się do tej samej godziny. Jest to
kluczowy element analizy syntaktycznej: program nie interpretuje jeszcze
znaczenia zjawisk pogodowych, lecz sprawdza zgodność dwóch uporządkowanych
ciągów liczbowych. W analizowanym przebiegu wszystkie źródła zostały
wyrównane do zakresu:
\[
2026\text{-}05\text{-}18\ 22{:}00
\quad \text{--} \quad
2026\text{-}05\text{-}20\ 23{:}00
\]
co daje 50 wspólnych godzin.

\noindent
Do oceny wykorzystano dwie miary. Pierwszą jest błąd RMSE:
\[
RMSE = \sqrt{\frac{1}{n}\sum_{i=1}^{n}(x_i - y_i)^2}
\]
gdzie $x_i$ oznacza wartość z modelu porównywanego, a $y_i$ wartość
referencyjną. Drugą miarą jest prosty system punktowy. Dla temperatury i wiatru
zastosowano progi:
\[
3\ pkt \leq 0.5,\qquad
2\ pkt \leq 1.5,\qquad
1\ pkt \leq 3.0
\]
natomiast dla opadów zastosowano ostrzejsze progi:
\[
3\ pkt \leq 0.1,\qquad
2\ pkt \leq 0.5,\qquad
1\ pkt \leq 2.0
\]
ponieważ opady są mierzone w milimetrach i zwykle mają mniejsze wartości
liczbowe niż temperatura lub prędkość wiatru.

\noindent
Poniższy fragment kodu pokazuje najważniejszą część porównywania. Nie jest to
cały program, ponieważ najistotniejsze są trzy operacje: wczytanie danych,
dopasowanie ich po czasie oraz obliczenie błędu i punktów.

\begin{lstlisting}[language=Python]
def przyznaj_punkty(blad):
    if blad <= 0.5:
        return 3
    if blad <= 1.5:
        return 2
    if blad <= 3.0:
        return 1
    return 0

def analyze_source(ground_truth, source_config):
    source = read_temperature(
        source_config["path"],
        source_config["key"]
    )

    merged = ground_truth.merge(source, on="time")

    source_column = f"temp_{source_config['key']}"
    merged["blad"] = (
        merged["temp_ground_truth"] - merged[source_column]
    ).abs()
    merged["punkty"] = merged["blad"].apply(przyznaj_punkty)

    points = merged["punkty"].sum()
    max_points = len(merged) * 3
    rmse = np.sqrt((merged["blad"] ** 2).mean())

    return points, max_points, rmse, len(merged)
\end{lstlisting}

\noindent
Funkcja \texttt{merge(..., on="time")} odpowiada za dopasowanie danych godzina
do godziny. Jeżeli jakieś źródło nie ma danej godziny, ta godzina nie jest
używana w rankingu. W aktualnej wersji skryptu dane zostały jednak wcześniej
przycięte do wspólnego zakresu, dlatego wszystkie modele mają po 50
porównywalnych obserwacji.

\subsubsection{Wykresy porównawcze}

\begin{figure}[!h]
\centering
\includegraphics[width=\textwidth]{wykres_temperatura.png}
\caption{Porównanie temperatury dla wszystkich analizowanych modeli.}
\end{figure}

\begin{figure}[!h]
\centering
\includegraphics[width=\textwidth]{wykres_wiatr.png}
\caption{Porównanie prędkości wiatru dla wszystkich analizowanych modeli.}
\end{figure}

\begin{figure}[!h]
\centering
\includegraphics[width=\textwidth]{wykres_opady.png}
\caption{Porównanie opadów dla wszystkich analizowanych modeli.}
\end{figure}

\clearpage
\subsubsection{Tabele rankingowe}

\begin{table}[!h]
\centering
\begin{tabular}{lrrrr}
\toprule
Model & Punkty & Wynik [\%] & RMSE [C] & Godziny \\
\midrule
Open-Meteo & 150/150 & 100.0 & 0.00 & 50 \\
GFS & 63/150 & 42.0 & 2.60 & 50 \\
AIFS & 46/150 & 30.7 & 3.05 & 50 \\
Yr.no & 43/150 & 28.7 & 3.33 & 50 \\
GraphCast & 7/150 & 4.7 & 7.25 & 50 \\
\bottomrule
\end{tabular}
\caption{Ranking modeli dla temperatury.}
\end{table}

\begin{table}[!h]
\centering
\begin{tabular}{lrrrr}
\toprule
Model & Punkty & Wynik [\%] & RMSE [m/s] & Godziny \\
\midrule
Open-Meteo & 150/150 & 100.0 & 0.00 & 50 \\
GFS & 133/150 & 88.7 & 0.60 & 50 \\
Yr.no & 125/150 & 83.3 & 0.76 & 50 \\
AIFS & 120/150 & 80.0 & 0.80 & 50 \\
GraphCast & 55/150 & 36.7 & 2.47 & 50 \\
\bottomrule
\end{tabular}
\caption{Ranking modeli dla prędkości wiatru.}
\end{table}

\begin{table}[!h]
\centering
\begin{tabular}{lrrrr}
\toprule
Model & Punkty & Wynik [\%] & RMSE [mm] & Godziny \\
\midrule
Open-Meteo & 150/150 & 100.0 & 0.00 & 50 \\
Yr.no & 150/150 & 100.0 & 0.00 & 50 \\
GFS & 150/150 & 100.0 & 0.02 & 50 \\
AIFS & 150/150 & 100.0 & 0.06 & 50 \\
GraphCast & 132/150 & 88.0 & 0.24 & 50 \\
\bottomrule
\end{tabular}
\caption{Ranking modeli dla opadów.}
\end{table}

\begin{table}[h!]
\centering
\small
\begin{tabular}{lrrp{6cm}}
\toprule
Model & Punkty łącznie & Wynik [\%] & Uwagi \\
\midrule
Open-Meteo & 450/450 & 100.0 & najlepsza zgodność z punktem odniesienia \\
GFS & 346/450 & 76.9 & dobre wyniki dla wiatru i opadów \\
Yr.no & 318/450 & 70.7 & dobre opady, słabsza temperatura \\
AIFS & 316/450 & 70.2 & dobre opady i wiatr, słabsza temperatura \\
GraphCast & 194/450 & 43.1 & najsłabszy wynik łączny \\
\bottomrule
\end{tabular}
\caption{Ranking łączny dla temperatury, wiatru i opadów.}
\end{table}

\noindent
Wynik Open-Meteo jest równy 100\%, ponieważ dane referencyjne
\texttt{default\_model} zostały utworzone na podstawie tego samego źródła.
Ten wynik należy traktować jako kontrolę poprawności działania programu:
jeżeli źródło porównywane jest identyczne z referencyjnym, błąd powinien być
równy zero. Dopiero pozostałe modele pokazują rzeczywiste różnice między
prognozami.

%--------------------------------------------------
\subsection{Interpretacja}

Otrzymane wyniki pokazują, że sama zgodność struktury danych nie gwarantuje
zgodności prognoz. Wszystkie modele zostały porównane dla tych samych godzin,
ale ich wyniki punktowe i RMSE różnią się w zależności od analizowanego
parametru.

\noindent
Najlepszy wynik uzyskało Open-Meteo, ponieważ katalog \texttt{default\_model}
został utworzony na podstawie tego samego źródła. Ten przypadek pełni więc
funkcję kontroli poprawności obliczeń. Jeżeli dane porównywane są identyczne
z danymi referencyjnymi, program powinien zwrócić maksymalną liczbę punktów
oraz RMSE równe zero.

\noindent
Wśród pozostałych modeli relatywnie dobrze wypadł GFS, szczególnie dla wiatru
i opadów. AIFS oraz yr.no uzyskały podobny wynik łączny, ale różniły się
rozkładem błędów między parametrami. Najsłabszy wynik łączny uzyskał
GraphCast, głównie przez duży błąd temperatury i wiatru.

\noindent
Opady okazały się parametrem, dla którego większość modeli uzyskała wysoką
punktację. Wynika to z niewielkich wartości opadów w analizowanym przedziale
czasowym. W takiej sytuacji nawet modele różniące się między sobą mogą
uzyskać zbliżone wyniki punktowe.

%--------------------------------------------------
\subsection{Wnioski}

Zastosowany sposób porównywania pozwala automatycznie zestawiać wiele źródeł
prognozy przy użyciu jednej procedury. Najważniejsze jest wcześniejsze
wyrównanie danych po czasie, ponieważ bez wspólnych godzin ranking nie miałby
sensu obliczeniowego.

\noindent
System punktowy jest prosty, ale czytelny. Pozwala szybko zbudować ranking,
natomiast RMSE daje dodatkową informację o skali błędu. Połączenie obu miar
sprawia, że wyniki można analizować zarówno jakościowo, jak i liczbowo.

\noindent
Projekt pokazuje, że dane z różnych serwisów i modeli AI można sprowadzić do
wspólnej reprezentacji syntaktycznej. Po takim ujednoliceniu dalsze
porównywanie nie zależy już od źródła danych, lecz od samych wartości
zapisanych w szeregach czasowych.


\section{Etap 3: Trenowanie własnego modelu prognostycznego}

\subsection{Cel budowy modelu}

Po przeprowadzeniu analizy porównawczej gotowych modeli prognostycznych kolejnym etapem projektu było opracowanie własnego modelu predykcyjnego wykorzystującego metody uczenia maszynowego. Celem tego etapu było sprawdzenie, czy model trenowany bezpośrednio na danych meteorologicznych jest w stanie poprawnie odwzorowywać krótkoterminowe zmiany parametrów atmosferycznych oraz konkurować z gotowymi rozwiązaniami prognostycznymi.

\noindent
Do realizacji tego zadania wykorzystano algorytm Random Forest Regressor dostępny w bibliotece \texttt{scikit-learn}. Random Forest należy do grupy metod zespołowych (\textit{ensemble learning}) i stanowi rozszerzenie klasycznych drzew decyzyjnych. Zamiast pojedynczego drzewa budowany jest cały zbiór niezależnych modeli, których wyniki są następnie uśredniane. Pozwala to znacząco ograniczyć problem przeuczenia (\textit{overfitting}) charakterystyczny dla pojedynczych drzew decyzyjnych.

\noindent
Algorytm Random Forest został wybrany ze względu na kilka istotnych właściwości. Przede wszystkim dobrze radzi sobie z modelowaniem zależności nieliniowych, które bardzo często występują w meteorologii. Parametry atmosferyczne, takie jak temperatura, wiatr czy opady, zależą jednocześnie od wielu czynników i wykazują skomplikowane relacje czasowe. Dodatkowo model ten nie wymaga silnych założeń matematycznych dotyczących rozkładu danych i dobrze radzi sobie nawet przy częściowo zaszumionych zbiorach.

\subsection{Przygotowanie danych wejściowych}

Proces budowy modelu rozpoczęto od przygotowania zbioru danych treningowych. Dane meteorologiczne zostały wcześniej pobrane z różnych źródeł prognostycznych oraz zsynchronizowane czasowo. Następnie wykonano etap oczyszczania danych obejmujący:

\begin{itemize}
    \item usuwanie brakujących rekordów,
    \item konwersję znaczników czasu do wspólnego formatu,
    \item sortowanie chronologiczne,
    \item eliminację zduplikowanych pomiarów,
    \item ujednolicenie nazw kolumn.
\end{itemize}

\noindent
Po oczyszczeniu danych przystąpiono do budowy zbioru cech wejściowych (\textit{features}). Model Random Forest nie analizuje danych w postaci surowego szeregu czasowego, dlatego konieczne było stworzenie dodatkowych zmiennych opisujących aktualny stan atmosfery oraz jego historię.

\noindent
W modelu wykorzystano następujące grupy cech:

\begin{itemize}
    \item aktualną temperaturę,
    \item prędkość wiatru,
    \item wielkość opadów,
    \item godzinę pomiaru,
    \item dzień tygodnia,
    \item opóźnione wartości parametrów meteorologicznych (\textit{lag features}).
\end{itemize}

\noindent
Szczególnie istotne były cechy czasowe. Temperatura i inne parametry meteorologiczne wykazują wyraźną cykliczność dobową, dlatego dodanie informacji o godzinie pomiaru umożliwia modelowi identyfikację regularnych wzorców zmian. Z kolei dzień tygodnia pozwala częściowo uwzględnić efekty sezonowe oraz różnice pomiędzy poszczególnymi dniami.

\noindent
Duże znaczenie miały również opóźnione wartości parametrów meteorologicznych. W praktyce oznacza to wykorzystanie danych z poprzednich godzin jako dodatkowych cech wejściowych. Jeżeli temperatura w ostatnich kilku godzinach rosła, model może wykorzystać tę informację do przewidywania dalszego trendu. Takie podejście pozwala częściowo odwzorować pamięć szeregu czasowego.

\noindent
Po przygotowaniu wszystkich cech dane podzielono na zbiór treningowy oraz testowy. Zbiór treningowy wykorzystywano do uczenia modelu, natomiast zbiór testowy służył do niezależnej oceny jakości predykcji. Rozdzielenie danych jest bardzo ważne, ponieważ pozwala sprawdzić, czy model potrafi poprawnie generalizować wiedzę na nowe, wcześniej niewidziane dane.

\subsection{Proces trenowania modelu}

Proces trenowania polegał na budowie wielu niezależnych drzew regresyjnych. Każde drzewo było trenowane na losowo wybranym podzbiorze danych treningowych oraz na losowym podzbiorze cech wejściowych. Mechanizm ten nazywany jest \textit{baggingiem} (\textit{Bootstrap Aggregating}) i stanowi podstawę działania Random Forest.

\noindent
Każde drzewo decyzyjne tworzy własny zestaw reguł podziału danych. Podczas budowy drzewa algorytm wybiera takie cechy i wartości progowe, które najlepiej redukują błąd predykcji. Przykładowo model może nauczyć się reguły wskazującej, że wysoka temperatura przy jednoczesnym niskim zachmurzeniu zwiększa prawdopodobieństwo dalszego wzrostu temperatury w kolejnych godzinach.

\noindent
Pojedyncze drzewa są stosunkowo niestabilne i podatne na przeuczenie, jednak połączenie dużej liczby drzew znacząco poprawia jakość predykcji. Wynik końcowy modelu regresyjnego wyznaczany jest jako średnia prognoz generowanych przez wszystkie drzewa należące do lasu.

\noindent
Podczas eksperymentów analizowano wpływ parametrów modelu na jakość wyników. Szczególną uwagę zwrócono na:

\begin{itemize}
    \item liczbę drzew decyzyjnych,
    \item maksymalną głębokość drzew,
    \item minimalną liczbę próbek wymaganych do podziału,
    \item poziom losowości procesu uczenia.
\end{itemize}

\noindent
Zwiększenie liczby drzew poprawiało stabilność predykcji oraz zmniejszało wariancję modelu, jednak jednocześnie wydłużało czas trenowania. Zbyt głębokie drzewa prowadziły natomiast do nadmiernego dopasowania modelu do danych treningowych.

\subsection{Ocena jakości predykcji}

Po zakończeniu trenowania model został poddany ewaluacji na zbiorze testowym. Jako podstawową miarę jakości wykorzystano wskaźnik RMSE (\textit{Root Mean Square Error}), który określa średnią skalę błędu prognozy. Metryka ta silniej penalizuje duże błędy, dzięki czemu dobrze nadaje się do oceny jakości modeli meteorologicznych.

\noindent
Oprócz analizy numerycznej przeprowadzono również wizualną ocenę jakości predykcji. W tym celu wygenerowano wykresy przedstawiające przebieg rzeczywistych wartości parametrów meteorologicznych oraz wartości przewidywanych przez model Random Forest.

\noindent
Analiza wykresów pozwoliła ocenić, czy model poprawnie odwzorowuje:

\begin{itemize}
    \item trendy dobowe,
    \item lokalne maksima i minima,
    \item amplitudę zmian temperatury,
    \item krótkookresowe fluktuacje parametrów atmosferycznych.
\end{itemize}

\noindent
Uzyskane wyniki wskazują, że model dobrze radzi sobie z odwzorowaniem ogólnego trendu zmian temperatury oraz częściowo poprawnie modeluje zmiany prędkości wiatru. Największe trudności pojawiały się podczas prognozowania opadów, które charakteryzują się wysoką nieregularnością oraz dużą zmiennością czasową.

\noindent
Pomimo tych ograniczeń Random Forest wykazał stosunkowo dobrą stabilność predykcji i niewielkie odchylenia od danych rzeczywistych. Wyniki te potwierdzają, że algorytmy uczenia maszynowego mogą stanowić skuteczne narzędzie wspomagające krótkoterminowe prognozowanie meteorologiczne.

\newpage

\section*{Podział pracy w zespole}

\begin{itemize}

    \item \textbf{Maciek Chotkowski} -- wsparcie w opracowaniu wyników oraz dopracowanie raportu

    \item \textbf{Martsin Lazouski} -- opracowanie procedur synchronizacji danych czasowych oraz implementacja obliczeń RMSE i systemu punktowej oceny modeli, modeli pozyskiwanie informacji (implemenctacja w Python) , przygotowanie struktury raportu.

    \item \textbf{Nadia Ziecina} -- wsparcie w opracowaniu wyników oraz dopracowanie raportu

    \item \textbf{Michał Kasznia} -- przygotowanie dokumentacji projektu

    \item \textbf{Julian Kasiński} -- ujednolicenie wyglądu raportu w tym szaty graficznej, rozwijanie koncepcji projektu

    \item \textbf{Volodymyr Kamenchuk} -- testowanie poprawności działania systemu, analiza wyników eksperymentów oraz walidacja danych wejściowych,implementacja analizy modeli AI (GraphCast, AIFS) oraz integracja danych pochodzących z różnych źródeł prognostycznych.

    \item \textbf{Usevalad Nikalaichuk} -- przygotowanie dokumentacji projektu, redakcja raportu oraz opracowanie części teoretycznej dotyczącej analizy danych meteorologicznych.

\end{itemize}

\end{document}
